\documentclass[11pt]{article}
\usepackage[margin=1in]{geometry}
\usepackage{amsmath,amssymb,amsthm,hyperref}
\begin{document}
\begin{center}{\Large Real rank zero does not mean projection density}\\[4pt]
Disproof of conjecture 00000000920\\[4pt]
ziangni-sys --- October 8, 2026\end{center}

Let $A=\mathbb C$ with its usual norm and involution $z^*=\overline z$.
This is an actual homogeneous $C^*$-algebra:
\[
 \mathbb C\cong C(\{\mathrm{pt}\},M_1(\mathbb C)),
 \qquad f\longmapsto f(\mathrm{pt}).
\]
Evaluation preserves multiplication, involution and norm, since the
supremum norm on a singleton equals the norm at that point. In particular,
$A$ is AH: it is the inductive limit of the constant homogeneous system
$A\xrightarrow{\mathrm{id}}A\xrightarrow{\mathrm{id}}\cdots$. No
infinite-dimensional restriction appears in the conjecture.

\paragraph{Real rank zero.}
The self-adjoint elements of $A$ are the real numbers. Given a self-adjoint
$a$ and $\varepsilon>0$, if $a\ne0$ take $b=a$. It is self-adjoint and
invertible, and $\|b-a\|=0<\varepsilon$. If $a=0$, take the real scalar
$b=\varepsilon/2$. It is nonzero, hence invertible, and
$\|b-a\|=\varepsilon/2<\varepsilon$. Thus every self-adjoint element is
approximable by self-adjoint invertibles, proving real rank zero in precisely
the approximation formulation stated in the source.

\paragraph{Projections are not dense.}
A projection $p\in\mathbb C$ satisfies $p^*=p$ and $p^2=p$. The latter
equation gives $p(p-1)=0$, so $p=0$ or $p=1$. Conversely both are
projections. For the self-adjoint element $a=1/2$,
\[
 \|p-a\|=\frac12\qquad\text{for every projection }p.
\]
Consequently the open norm ball of radius $1/4$ around $1/2$ contains no
projection. The projection semilattice is exactly $\{0,1\}$, with its usual
order, and is therefore not dense among self-adjoint elements.

This disproves the claimed equivalence: the AH algebra has real rank zero
but fails the stated density condition. The further cancellation description
cannot restore an already-false equivalence.

\paragraph{Interpretation of the density clause.}
The Chinese version explicitly uses the projection semilattice, while the
English version calls it the projection semigroup. Neither notion means
arbitrary real linear combinations of projections. Even the more generous
interpretation as finite additive sums of projections fails here: such sums
are nonnegative integers, each at distance at least $1/2$ from $1/2$.
The stable Murray--von Neumann semigroup $V(A)$ is an abstract semigroup
of classes, not itself a normed subset of $A$. The counterexample addresses
the literal projection-density assertion in the source. The familiar
real-rank-zero characterization by approximation with self-adjoint
\emph{finite-spectrum} elements, or by suitable real linear combinations
of projections, is a different condition and is not refuted.

\paragraph{Lean certificate.}
The project constructs the actual continuous-function algebra on a point,
its algebra equivalence with $\mathbb C$ and the universal property of the
constant identity system. It proves self-adjoint invertible approximation,
classifies all complex projections, computes the exact norm gap at $1/2$
and proves real rank zero together with failure of projection density.
\end{document}
